\documentclass[../../main.tex]{subfiles}
\begin{document}

\begin{center}
  {\fontsize{48pt}{58pt}\selectfont \textbf{2012}\normalsize\textbf{年度}}
\end{center}

\begin{table}[htb]
  \centering
  \caption{出題テーマと難易度}
  \begin{tabular}{|c|p{9cm}|c|}
    \hline
    \textbf{問題番号} & \textbf{テーマ} & \textbf{難易度} \\
    \hline
    第1問 & 冪根の極限と対数を含む定積分 & \\
    \hline
    第2問 & 正四面体内の正三角形と平行線 & \\
    \hline
    第3問 & 二次式の条件下での三次式の値域 & \\
    \hline
    第4問 & 三乗根の無理性と多項式の整除 & \\
    \hline
    第5問 & 正多角形と三角形の命題の真偽 & \\
    \hline
    第6問 & サイコロで定まる連分数型の確率 & \\
    \hline
  \end{tabular}
\end{table}

\begin{table}[htb]
  \centering
  \caption{解答}
  \renewcommand{\arraystretch}{1.4}
  \begin{tabular}{|c|c|p{9cm}|}
    \hline
    \textbf{問題番号} & \textbf{小問} & \textbf{答え} \\
    \hline
    第1問 & (1) & $\max(1,a)$ \\
    \cline{2-3}
    & (2) & $\dfrac{\pi}{12}+\dfrac{1}{2}\left(1-\dfrac{2\sqrt{3}}{3}\right)\log2$ \\
    \hline
    第2問 & -- & $PQ\parallel AB$，$QR\parallel BC$，$RP\parallel CA$（証明問題） \\
    \hline
    第3問 & -- & $-2(3\sqrt{2}+4)\le x^2y+xy^2-x^2-2xy-y^2+x+y\le3$ \\
    \hline
    第4問 & (1) & $\sqrt[3]{2}$は無理数（証明問題） \\
    \cline{2-3}
    & (2) & $P(x)$は$x^3-2$で割り切れる（証明問題） \\
    \hline
    第5問 & (p) & 正しい \\
    \cline{2-3}
    & (q) & 正しくない（正六角形の頂点による反例） \\
    \hline
    第6問 & -- & $p_n=\dfrac{1}{5}\left(1-6^{-n}\right)$ \\
    \hline
  \end{tabular}
\end{table}

\paragraph{試験形態}
6問．試験時間は要確認．

\end{document}
